(1) 半径 \(r\) 的竖直细管自液面向上吸入润湿液体，表面张力 \(\gamma\)、接触角 \(0\le\theta<\pi/2\)、动力黏度 \(\mu\)、密度 \(\rho\)。忽略入口阻力、惯性和气相黏性，取浸润高度 \(L(t)\)。由Laplace压差、重力静压及Poiseuille流阻推导移动前沿的非线性微分方程。
(2) 先令 \(g=0\)、\(L(0)=0\)，精确求 \(L(t)\) 的平方根律及前沿速度，解释 \(t=0\) 速度发散的近似来源并核验量纲。
(3) 对 \(g>0\) 求静态极限高度 \(L_\infty\)，将时间积分为高度的隐式闭式；证明 \(0<L<L_\infty\) 时单调上升并求接近平衡的指数时间常数。
(4) 令体积流率 \(Q=\pi r^2\dot L\)。从流体压力功与黏性损失独立验证毛细功率等于重力势能增长率加黏性耗散率，说明在何种尺度必须恢复液柱惯性或动态接触角。
